\documentclass[11pt]{article}

\usepackage[a4paper,margin=1in]{geometry}
\usepackage{amsmath,amssymb,amsthm,mathtools}
\usepackage{booktabs}
\usepackage{array}
\usepackage[hidelinks]{hyperref}
\usepackage{microtype}

\newtheorem{theorem}{Theorem}
\newtheorem{lemma}{Lemma}
\newtheorem{corollary}{Corollary}
\theoremstyle{remark}
\newtheorem{remark}{Remark}

\DeclareMathOperator{\LHS}{LHS}

\title{An Infinite Family of Integer Solutions to \texorpdfstring{\(z^2+y^2z+2x^3+1=0\)}{z^2+y^2z+2x^3+1=0}}
\author{}
\date{}

\begin{document}
\maketitle

\begin{abstract}
We give a fully explicit infinite family of integer solutions to
\[
        z^2+y^2z+2x^3+1=0.
\]
The proof is elementary and self-contained.  The construction is based on a polynomial identity which makes
\(2X^3-1\) split as
\((M-2T^2)(M+2T^2)\).  Setting \(x=-X\), \(y=2T\), and
\(z=M-2T^2\) then gives the required equation.
\end{abstract}

\section{Statement of the result}

The goal is not to classify all integer solutions, but to prove unconditionally that infinitely many exist.

\begin{theorem}\label{thm:main}
The Diophantine equation
\[
        z^2+y^2z+2x^3+1=0
\]
has infinitely many integer solutions.  More precisely, for every odd integer \(r\geq 3\), define
\[
        a=r^2-2,\qquad n=\frac{a^2-1}{2},\qquad T=\frac{r(a-1)}2,
\]
\[
        X=2n^2+a,
        \qquad
        M=4n^3+3an+1.
\]
Then
\[
        \boxed{(x,y,z)=\bigl(-X,\,2T,\,M-2T^2\bigr)}
\]
is an integer solution of
\[
        z^2+y^2z+2x^3+1=0.
\]
As \(r\) ranges over the odd integers \(r\geq 3\), these solutions are pairwise distinct.
\end{theorem}

The same family can be written directly as a polynomial-parametric family:
\begin{corollary}\label{cor:polynomial}
For every odd integer \(r\geq 3\), the triple
\[
\boxed{
\begin{aligned}
 x_r&=-\frac{r^8-8r^6+22r^4-22r^2+5}{2},\\
 y_r&=r^3-3r,\\
 z_r&=\frac{r^{12}-12r^{10}+57r^8-134r^6+159r^4-84r^2+11}{2}
\end{aligned}}
\]
is an integer solution of
\[
        z^2+y^2z+2x^3+1=0.
\]
\end{corollary}

\section{The auxiliary identity}

We first isolate the identity behind the construction.

\begin{lemma}\label{lem:identity}
Let \(a\) and \(n\) be rational numbers satisfying
\[
        n=\frac{a^2-1}{2}.
\]
Define
\[
        X=2n^2+a,
        \qquad
        M=4n^3+3an+1.
\]
Then
\[
        M^2-2X^3+1=\frac{(a-1)^4(a+2)^2}{4}.
\]
\end{lemma}

\begin{proof}
We expand first without substituting the special value of \(n\).  Since
\[
        M=4n^3+3an+1,
        \qquad
        X=2n^2+a,
\]
we have
\begin{align*}
 M^2
 &= (4n^3+3an+1)^2 \\
 &=16n^6+24an^4+8n^3+9a^2n^2+6an+1,
\end{align*}
and
\begin{align*}
 2X^3
 &=2(2n^2+a)^3 \\
 &=2(8n^6+12an^4+6a^2n^2+a^3) \\
 &=16n^6+24an^4+12a^2n^2+2a^3.
\end{align*}
Therefore
\begin{align}\label{eq:basic-expanded}
 M^2-2X^3+1
 &=8n^3-3a^2n^2+6an-2a^3+2.
\end{align}
Now substitute \(n=(a^2-1)/2\).  Then \eqref{eq:basic-expanded} becomes
\begin{align*}
&8\left(\frac{a^2-1}{2}\right)^3
 -3a^2\left(\frac{a^2-1}{2}\right)^2
 +6a\left(\frac{a^2-1}{2}\right)
 -2a^3+2 \\
&\qquad = (a^2-1)^3-\frac{3a^2(a^2-1)^2}{4}+3a(a^2-1)-2a^3+2.
\end{align*}
Multiplying by \(4\), the last expression is equal to
\begin{align*}
&4(a^2-1)^3-3a^2(a^2-1)^2+12a(a^2-1)-8a^3+8 \\
&\qquad = a^6-6a^4+4a^3+9a^2-12a+4 \\
&\qquad = (a-1)^4(a+2)^2.
\end{align*}
Dividing by \(4\) gives
\[
        M^2-2X^3+1=\frac{(a-1)^4(a+2)^2}{4},
\]
as claimed.
\end{proof}

\section{Turning the identity into a solution}

The next lemma is the conversion step.  It shows how a factorization of \(2X^3-1\) immediately gives a solution of the original equation.

\begin{lemma}\label{lem:conversion}
Let \(X,M,T\in\mathbb Z\) satisfy
\[
        M^2-2X^3+1=4T^4.
\]
Then
\[
        (x,y,z)=(-X,\,2T,\,M-2T^2)
\]
is an integer solution of
\[
        z^2+y^2z+2x^3+1=0.
\]
\end{lemma}

\begin{proof}
Define
\[
        x=-X,
        \qquad
        y=2T,
        \qquad
        z=M-2T^2.
\]
Then
\[
        z+y^2=M-2T^2+(2T)^2=M+2T^2.
\]
Therefore
\begin{align*}
        z(z+y^2)
        &=(M-2T^2)(M+2T^2) \\
        &=M^2-4T^4.
\end{align*}
Using the assumed identity \(M^2-2X^3+1=4T^4\), we get
\[
        M^2-4T^4=2X^3-1.
\]
Thus
\begin{align*}
 z^2+y^2z+2x^3+1
 &=z(z+y^2)+2(-X)^3+1 \\
 &=(2X^3-1)-2X^3+1 \\
 &=0.
\end{align*}
This proves the lemma.
\end{proof}

\section{Proof of the infinite family}

\begin{proof}[Proof of Theorem \ref{thm:main}]
Let \(r\geq 3\) be odd.  Define
\[
        a=r^2-2,
        \qquad
        n=\frac{a^2-1}{2},
        \qquad
        T=\frac{r(a-1)}2,
\]
\[
        X=2n^2+a,
        \qquad
        M=4n^3+3an+1.
\]
We first check that all quantities are integral.  Since \(r\) is odd, \(r^2\) is odd, so
\[
        a=r^2-2
\]
is odd.  Hence \(a^2-1=(a-1)(a+1)\) is divisible by \(8\), and in particular
\[
        n=\frac{a^2-1}{2}\in\mathbb Z.
\]
Also \(a-1=r^2-3\) is even, so
\[
        T=\frac{r(a-1)}2\in\mathbb Z.
\]
It follows immediately that \(X,M\in\mathbb Z\).

By Lemma \ref{lem:identity},
\[
        M^2-2X^3+1=\frac{(a-1)^4(a+2)^2}{4}.
\]
Because \(a=r^2-2\), we have \(a+2=r^2\).  Hence
\begin{align*}
        \frac{(a-1)^4(a+2)^2}{4}
        &=\frac{(a-1)^4r^4}{4} \\
        &=4\left(\frac{r(a-1)}2\right)^4 \\
        &=4T^4.
\end{align*}
Thus
\[
        M^2-2X^3+1=4T^4.
\]
Lemma \ref{lem:conversion} now gives that
\[
        (x,y,z)=(-X,2T,M-2T^2)
\]
is an integer solution of
\[
        z^2+y^2z+2x^3+1=0.
\]

It remains only to prove that the construction gives infinitely many distinct solutions.  For \(r\geq 3\), the number
\[
        a=r^2-2
\]
is strictly increasing with \(r\), and
\[
        n=\frac{a^2-1}{2}
\]
is strictly increasing with \(a\) for \(a\geq 7\).  Therefore
\[
        X=2n^2+a
\]
is strictly increasing as \(r\) ranges over the odd integers \(r\geq 3\).  Consequently the first coordinate
\[
        x=-X
\]
assumes infinitely many distinct values.  The corresponding triples are therefore pairwise distinct.
\end{proof}

\begin{proof}[Proof of Corollary \ref{cor:polynomial}]
Substituting \(a=r^2-2\) into
\[
        n=\frac{a^2-1}{2}
\]
yields
\[
        n=\frac{(r^2-2)^2-1}{2}
         =\frac{(r^2-3)(r^2-1)}{2}.
\]
Then
\begin{align*}
        X
        &=2n^2+a \\
        &=2\left(\frac{(r^2-3)(r^2-1)}{2}\right)^2+r^2-2 \\
        &=\frac{r^8-8r^6+22r^4-22r^2+5}{2}.
\end{align*}
Also
\[
        2T=r(a-1)=r(r^2-3)=r^3-3r.
\]
Finally,
\begin{align*}
        M-2T^2
        &=4\left(\frac{(r^2-3)(r^2-1)}{2}\right)^3
          +3(r^2-2)\left(\frac{(r^2-3)(r^2-1)}{2}\right)+1
          -2\left(\frac{r(r^2-3)}{2}\right)^2 \\
        &=\frac{r^{12}-12r^{10}+57r^8-134r^6+159r^4-84r^2+11}{2}.
\end{align*}
These expressions are exactly the stated \((x_r,y_r,z_r)\).  By Theorem \ref{thm:main}, they solve the equation.
\end{proof}

\section{Examples}

The following table lists several integer solutions produced by the construction.  The final column records the value of
\[
        z^2+y^2z+2x^3+1,
\]
which is \(0\) in each case.

\begin{center}
\begin{tabular}{@{}r r r r r@{}}
\toprule
\(r\) & \(x\) & \(y\) & \(z\) & \(z^2+y^2z+2x^3+1\) \\
\midrule
\(3\)  & \(-1159\)       & \(18\)   & \(55639\)          & \(0\) \\
\(5\)  & \(-139415\)     & \(110\)  & \(73611143\)       & \(0\) \\
\(7\)  & \(-2437679\)    & \(322\)  & \(5382395279\)     & \(0\) \\
\(9\)  & \(-19468879\)   & \(702\)  & \(121485805039\)   & \(0\) \\
\(11\) & \(-100252919\)  & \(1298\) & \(1419581333159\)  & \(0\) \\
\bottomrule
\end{tabular}
\end{center}

For instance, the first row gives
\[
        (x,y,z)=(-1159,18,55639).
\]
A direct check may be written in the factorized form
\[
        55639(55639+18^2)=3113725357=2\cdot 1159^3-1.
\]
Therefore
\begin{align*}
        55639^2+18^2\cdot 55639+2(-1159)^3+1
        &=\bigl(2\cdot 1159^3-1\bigr)-2\cdot1159^3+1 \\
        &=0.
\end{align*}

\section{A sign symmetry}

Since the original equation contains \(y\) only through \(y^2\), every solution \((x,y,z)\) immediately gives another solution \((x,-y,z)\).  Thus the displayed family also yields the companion solutions
\[
        (-X,-2T,M-2T^2)
\]
for the same odd parameters \(r\geq 3\).

\section{Conclusion}

The elementary identity in Lemma \ref{lem:identity} produces infinitely many factorizations
\[
        2X^3-1=(M-2T^2)(M+2T^2).
\]
The substitution
\[
        x=-X,
        \qquad
        y=2T,
        \qquad
        z=M-2T^2
\]
then converts these factorizations into integer solutions of
\[
        z^2+y^2z+2x^3+1=0.
\]
Because the parameter \(r\) can be chosen as any odd integer \(r\geq 3\), and because the corresponding \(x\)-coordinates are pairwise distinct, the equation has infinitely many integer solutions.

\end{document}
